\chapter{Fields, Particles, and the Klein--Gordon Equation}

This lecture begins in an unfinished place and ends with the outline of a relativistic field theory. A student asks about the Doppler shift, and Susskind turns that question into a quick exercise in Lorentz geometry: do not memorize the formula if the spacetime picture can generate it for us. Only then does he reset the room and begin the advertised subject, \emph{fields and particles}. From there the lecture unfolds in a very deliberate sequence: a field acts on a particle, the particle must act back on the field, the common action tells us exactly how, the notation of relativity is introduced when the component formulas become too cumbersome, and a small extension of the scalar-field Lagrangian produces the Klein--Gordon equation. These notes follow Leonard Susskind's lecture closely, with the transcript as the primary guide and with selected board frames preserved from the lecture record by LazyingArt LLC and Video2Book where they genuinely support the mathematics.

\section{Relativistic Doppler Shift from Lorentz Geometry}

The lecture opens with a correction to a vague memory: a fast observer does \emph{not} always see a longer wavelength. It depends on the direction of motion. If we move with the light ray, the wavelength gets stretched; if we move against it, it gets compressed. Susskind does not present this as a formula to be recalled. He presents it as a geometric fact about different notions of simultaneity.

We take two inertial frames. The unprimed frame is at rest. The primed frame moves to the right with speed \(V\), and for the moment we use units in which \(c=1\). A right-moving light ray follows the line
\begin{equation}
X=T.
\end{equation}
A wave is represented by a family of such lightlike crest lines. If the wavelength in the unprimed frame is normalized to \(1\), then two neighboring crests may be written as
\begin{equation}
X=T,
\qquad
X=T+1.
\end{equation}
In the unprimed frame the wavelength is the spatial distance between these crest lines at fixed \(T\). In the moving frame it is the spatial distance between them at fixed \(T'\). The whole calculation is therefore a calculation about simultaneity.

\subsection{Question \& Answer}

How do we compute the wavelength in the moving frame without guessing?

We use the Lorentz transformation directly:
\begin{equation}
T'=\frac{T-VX}{\sqrt{1-V^2}},
\qquad
X'=\frac{X-VT}{\sqrt{1-V^2}}.
\end{equation}
An instant of time in the primed frame is described by \(T'=\text{constant}\). Taking the convenient slice \(T'=0\), we obtain
\begin{equation}
T=VX.
\end{equation}
This is the line of primed simultaneity we need.

Now take the neighboring crest \(X=T+1\) and intersect it with \(T=VX\). Then
\begin{align}
X &= VX+1,\\
X(1-V) &= 1,\\
X &= \frac{1}{1-V}.
\end{align}
The corresponding time coordinate is therefore
\begin{equation}
T=VX=\frac{V}{1-V}.
\end{equation}
At that event the primed spatial coordinate is
\begin{align}
X'
&=
\frac{X-VT}{\sqrt{1-V^2}}\\
&=
\frac{\frac{1}{1-V}-V\frac{V}{1-V}}{\sqrt{1-V^2}}\\
&=
\frac{1-V^2}{(1-V)\sqrt{1-V^2}}\\
&=
\sqrt{\frac{1+V}{1-V}}.
\end{align}
Because we normalized the original wavelength to \(1\), this is already the wavelength ratio:
\begin{equation}
\frac{\lambda'}{\lambda}=\sqrt{\frac{1+V}{1-V}}.
\end{equation}
Restoring units by writing \(V=v/c\), we get the standard relativistic Doppler formula
\begin{equation}
\frac{\lambda'}{\lambda}
=
\sqrt{\frac{1+v/c}{1-v/c}}
=
\sqrt{\frac{c+v}{c-v}}.
\end{equation}

If we reverse the observer's direction, we simply replace \(v\) by \(-v\). The ratio then becomes less than \(1\): motion against the light ray shortens the observed wavelength. Motion with the light ray lengthens it, and as \(v\to c\) the observed wavelength diverges.

This is the lecture's opening lesson in miniature: it is all there in the picture, provided we identify the right events and translate the picture into equations.

After a short administrative interruption, Susskind wipes the slate clean and begins the main lecture.

\section{Fields, Particles, and Action-Reaction}

The new opening is abrupt: \emph{fields and particles}. In quantum mechanics the distinction eventually breaks down, but not here. Here a field fills spacetime and a particle traces a worldline. They are not the same thing, but they will meet in the action.

The question is posed in a deliberately elementary way. Suppose a field influences a particle, perhaps by exerting forces on it. Why must the particle influence the field in return? Why can the influence not run only one way?

\subsection{Question \& Answer}

Why does the action principle force influence to be mutual?

Susskind answers this first in a finite-dimensional toy model. Let the system have two coordinates \(x(t)\) and \(y(t)\), not necessarily spatial coordinates in the ordinary sense, just two generalized coordinates. Suppose the Lagrangian splits:
\begin{equation}
L(x,\dot x,y,\dot y)=L_x(x,\dot x)+L_y(y,\dot y).
\end{equation}
Then the Euler--Lagrange equations are
\begin{equation}
\frac{d}{dt}\frac{\partial L}{\partial \dot x}=\frac{\partial L}{\partial x},
\qquad
\frac{d}{dt}\frac{\partial L}{\partial \dot y}=\frac{\partial L}{\partial y}.
\end{equation}
But because \(L\) is a sum of separate \(x\)- and \(y\)-pieces, the \(x\)-equation contains no \(y\), and the \(y\)-equation contains no \(x\). The two degrees of freedom do not affect one another.

So if \(y\) does affect \(x\), then the Lagrangian must be more complicated. For example,
\begin{equation}
L=\frac{\dot x^2}{2}+\frac{\dot y^2}{2}-V(x)-V'(y)+xy.
\end{equation}
Now the mixed term \(xy\) appears in both Euler--Lagrange equations. The equation for \(x\) depends on \(y\), and the equation for \(y\) depends on \(x\). A common term in the Lagrangian creates common influence. That is action and reaction in its simplest Lagrangian form.

This toy model is not a digression. It is the finite-dimensional version of the field-particle story to come.

\section{A Particle in a Prescribed Scalar Field}

Susskind now returns to a scalar field \(\phi(x,t)\), taken at first as something already known. Somebody has solved its equations, or measured it, and the result is some function of position and time. We are not yet asking where the field comes from. We ask only: if a particle moves in that field, what Lagrangian should describe the motion?

The starting point is the familiar relativistic particle action
\begin{equation}
S_{\mathrm{particle}}=\int dt\,\mathcal L_{\mathrm{particle}}=-m\int d\tau,
\end{equation}
with free-particle Lagrangian
\begin{equation}
\mathcal L_{\mathrm{particle}}=-m\sqrt{1-\dot x^{\,2}}
\end{equation}
in units where \(c=1\).

\begin{figure}
\centering
\includegraphics[width=0.82\textwidth]{lecture_05_figure_02.png}
\end{figure}

This board frame is worth keeping because it shows exactly that free relativistic particle term. The left-hand label is partially hidden, but the square-root form is unambiguous. The little slanted sketch at the right is preserved as board context only; it is too stripped down to support a separate reconstruction.

To let the field influence the particle, Susskind proposes the simplest modification of the same invariant structure. The spoken blackboard wording is compressed, but the later nonrelativistic expansion fixes the sign:
\begin{equation}
\mathcal L_{\mathrm{particle}}
=
-(m+g\phi(x,t))\sqrt{1-\dot x^{\,2}}
\end{equation}
in \(c=1\) units. Restoring \(c\), the cleaner form is
\begin{equation}
\mathcal L_{\mathrm{particle}}
=
-(mc^2+g\phi(x,t))\sqrt{1-\frac{\dot x^{\,2}}{c^2}}.
\end{equation}
Here \(g\) is the coupling constant that measures the strength with which the field affects the particle.

The lecture then makes a very characteristic simplifying move. We go to one spatial dimension and take the nonrelativistic limit, not because that is the final theory, but because it exposes the meaning of the interaction. Expand the square root:
\begin{equation}
\sqrt{1-\frac{\dot x^{\,2}}{c^2}}
\approx
1-\frac{\dot x^{\,2}}{2c^2}.
\end{equation}
Substituting this into the particle Lagrangian gives
\begin{align}
\mathcal L_{\mathrm{particle}}
&\approx
-(mc^2+g\phi)\left(1-\frac{\dot x^{\,2}}{2c^2}\right)\\
&=
-mc^2-g\phi+\frac{m\dot x^{\,2}}{2}
+\mathcal O(c^{-2}).
\end{align}
The first term, \(-mc^2\), is just a constant. It does not affect the equations of motion and may be discarded. What remains is
\begin{equation}
L_{\mathrm{NR}}=\frac12 m\dot x^{\,2}-g\phi(x,t).
\end{equation}
Now the interpretation is immediate. This has the standard nonrelativistic form \(T-V\), so
\begin{equation}
V_{\mathrm{particle}}(x,t)=g\phi(x,t).
\end{equation}
Thus \(g\phi\) is the potential energy of the particle in the scalar field.

The lecture pauses briefly over units here. In this simplified setting \(g\phi\) must have units of energy. Susskind remarks that \(\phi\) therefore carries units of energy in the present example, and \(g\) is dimensionless here, though in other spacetime dimensions the bookkeeping can change.

So the first half of the story is clear. The field acts on the particle because the particle Lagrangian contains \(\phi\).

\section{The Particle as a Source of the Field}

Now comes the reversal. What about the effect of the particle on the field? The key sentence is: there is only one action. Not one action for the particle and another unrelated one for the field, but a single action for the combined system. To find the real motion we vary both the particle trajectory and the field configuration until the common action is stationary.

For the field, we begin with the scalar-field action written as an integral over spacetime:
\begin{equation}
S_f
=
\int d^4x\,
\left[
\frac12\left(\frac{\partial\phi}{\partial t}\right)^2
-\frac12\left(\frac{\partial\phi}{\partial x}\right)^2
-\frac12\left(\frac{\partial\phi}{\partial y}\right)^2
-\frac12\left(\frac{\partial\phi}{\partial z}\right)^2
\right].
\end{equation}
In the lecture this is first recalled in a reduced form and then extended verbally to all three spatial directions. Here we write the cleaned \(3+1\)-dimensional version explicitly.

To make the basic point without additional clutter, Susskind takes the particle to be at rest at the origin. Then the interaction part of the particle action is
\begin{equation}
S_{\mathrm{int}}=-g\int dt\,\phi(0,t).
\end{equation}
This is still an integral only along the particle worldline. But the field action is an integral over all spacetime. The question is then posed exactly as in the lecture: how do we turn the worldline term into a spacetime integral so that the whole thing can be written as one action?

The answer is the Dirac delta function. We write
\begin{equation}
\phi(0,t)=\int d^3x\,\phi(\mathbf x,t)\,\delta^3(\mathbf x),
\end{equation}
where
\begin{equation}
\delta^3(\mathbf x)=\delta(x)\delta(y)\delta(z),
\qquad
\int dx\,f(x)\delta(x)=f(0).
\end{equation}
This is the entire trick. The delta function simply picks out the value of the field at the particle's position. Therefore
\begin{equation}
S_{\mathrm{int}}=-g\int d^4x\,\delta^3(\mathbf x)\phi(\mathbf x,t).
\end{equation}
Now the full action has one common form,
\begin{equation}
S=S_f+S_{\mathrm{int}},
\end{equation}
and the same term \(g\phi\) that told the particle how to move will now tell the field how the particle acts as a source.

\begin{figure}
\centering
\includegraphics[width=0.82\textwidth]{lecture_05_figure_03.png}
\end{figure}

This board frame is useful precisely because it catches the live derivation in midstream. The right-hand side of the board already contains the finished second-derivative stack, while the left edge shows the next symbol just beginning to be written. The reliable mathematical content is the operator
\begin{equation}
\frac{\partial^2\phi}{\partial t^2}
-\frac{\partial^2\phi}{\partial x^2}
-\frac{\partial^2\phi}{\partial y^2}
-\frac{\partial^2\phi}{\partial z^2}.
\end{equation}

Varying the action with respect to \(\phi\), the derivative terms give the wave operator, while the interaction term contributes the source. With the lecture's corrected sign convention, the cleaned equation is
\begin{equation}
\frac{\partial^2\phi}{\partial t^2}-\nabla^2\phi=-g\,\delta^3(\mathbf x).
\end{equation}
This is a wave equation everywhere except at the location of the particle, where the delta function marks a special point in space.

To make the structure more familiar, Susskind immediately asks for a static solution. If the particle is standing still, we may look for a field that does not vary with time. Then \(\partial_t\phi=0\), so the time-derivative term drops out and the equation becomes
\begin{equation}
\nabla^2\phi=g\,\delta^3(\mathbf x).
\end{equation}
This is Poisson's equation with a point source. The analogy with electrostatics is deliberate: we are not yet doing electrodynamics, but we are seeing exactly how a source term appears in a field equation.

And now the reciprocity is explicit. The very same term in the action that makes the particle move as if it feels the potential energy \(g\phi\) also makes the field satisfy an equation with source \(-g\,\delta^3(\mathbf x)\). The same coupling constant \(g\) measures both effects.

\section{Moving Sources and Time Dependence}

Once the static example has done its job, the lecture loosens it slightly. Suppose the particle is not at the origin but fixed at \(x=a\). Then the one-dimensional source is shifted by the obvious replacement
\begin{equation}
\delta(x)\longrightarrow \delta(x-a).
\end{equation}
If the particle moves, with position \(x=a(t)\), then the source becomes
\begin{equation}
\delta(x)\longrightarrow \delta(x-a(t)).
\end{equation}
This change is small notationally, but physically it matters a great deal. A time-dependent source cannot be paired with a time-independent field. If the right-hand side of the field equation depends on time, then the left-hand side must carry that time dependence too. So the term
\begin{equation}
\frac{\partial^2\phi}{\partial t^2}
\end{equation}
must be restored. A moving, accelerating, or vibrating particle makes the field time dependent, and a time-dependent field radiates waves.

At this point the lecture refuses to solve the radiation problem. That is not evasion. It is a structural turn. The component formulas have become cumbersome, and this is exactly when Susskind decides it is time to stop writing every index out by hand and introduce the notation of relativity in a serious way.

\section{Relativistic Notation as a Working Tool}

The lecture is very explicit about why this section appears here. We are not introducing notation because notation is glamorous. We are introducing it because we are tired of writing the second derivative with respect to \(x\), then with respect to \(y\), then with respect to \(z\), every single time we touch a field equation. The right notation will compress the formulas and at the same time keep Lorentz invariance visible.

We denote the spacetime coordinates by
\begin{equation}
X^\mu=(t,x,y,z),
\end{equation}
with Greek indices \(\mu,\nu\) running over all four coordinates. A four-vector is any four-component object that transforms under Lorentz transformations in the same way as \(X^\mu\).

Out of a four-vector we can build a scalar. The standard example is proper time:
\begin{equation}
d\tau^2=dt^2-dx^2-dy^2-dz^2.
\end{equation}
To write such expressions neatly, we introduce the metric
\begin{equation}
\eta_{\mu\nu}=\mathrm{diag}(-1,1,1,1).
\end{equation}
It acts on a contravariant vector \(A^\mu\) to produce a covariant one:
\begin{equation}
A_\mu=\eta_{\mu\nu}A^\nu.
\end{equation}
With this signature convention, lowering the index changes the sign of the time component and leaves the spatial components alone.

\begin{figure}
\centering
\includegraphics[width=0.82\textwidth]{lecture_05_figure_04.png}
\end{figure}

The board frame here is especially valuable because it shows the two pieces of the notation lesson at once. On the left is the contraction that makes a scalar. On the right is the differential of a scalar field written as its derivative paired with \(dx^\mu\). The clean version of the contraction is
\begin{equation}
A^\mu A_\mu=A^\mu\eta_{\mu\nu}A^\nu.
\end{equation}
This is the Minkowski analog of an ordinary dot product, with the metric supplying the crucial minus sign.

The second notational device is Einstein summation: whenever the same index appears once upstairs and once downstairs, we sum over it. This lets us write contractions without cluttering the page with explicit summation symbols.

Now take a scalar field \(\phi(X)\). The infinitesimal change in \(\phi\) between neighboring points is
\begin{equation}
d\phi=\frac{\partial\phi(X)}{\partial x^\mu}\,dx^\mu.
\end{equation}
Because \(d\phi\) is a scalar and \(dx^\mu\) is a contravariant four-vector, the coefficient of \(dx^\mu\) must transform as a covariant four-vector. The lecture states the underlying theorem in words: if \(a^\mu b_\mu\) is a scalar for every four-vector \(a^\mu\), then \(b_\mu\) must itself be a covariant four-vector. We apply that theorem here to conclude that
\begin{equation}
\partial_\mu\phi=\frac{\partial\phi}{\partial x^\mu}
=
\left(
\frac{\partial\phi}{\partial t},
\frac{\partial\phi}{\partial x},
\frac{\partial\phi}{\partial y},
\frac{\partial\phi}{\partial z}
\right)
\end{equation}
is a covariant vector. Raising the index gives
\begin{equation}
\partial^\mu\phi
=
\left(
-\frac{\partial\phi}{\partial t},
\frac{\partial\phi}{\partial x},
\frac{\partial\phi}{\partial y},
\frac{\partial\phi}{\partial z}
\right).
\end{equation}
Therefore
\begin{equation}
\partial_\mu\phi\,\partial^\mu\phi
=
-\left(\frac{\partial\phi}{\partial t}\right)^2
+\left(\frac{\partial\phi}{\partial x}\right)^2
+\left(\frac{\partial\phi}{\partial y}\right)^2
+\left(\frac{\partial\phi}{\partial z}\right)^2
\end{equation}
is a scalar, and the scalar-field Lagrangian density may be written compactly as
\begin{equation}
\mathcal L=-\frac12\,\partial_\mu\phi\,\partial^\mu\phi.
\end{equation}
This is exactly the same content as the long component expression, but now the Lorentz structure is visible at a glance.

\subsection{Question \& Answer}

Why should the Lagrangian be a scalar?

Because if the Lagrangian density is a scalar, then the action obtained by integrating it over spacetime is the same in every inertial frame. And if the action is frame-independent, then the principle of least action produces equations of motion with the same form in every frame. In this way the invariance of the Lagrangian is what enforces the invariance of the laws.

The lecture inserts a brief but worthwhile aside here. Newtonian mechanics also has a relativity principle. If we take the limit \(c\to\infty\) in the Lorentz transformations, we recover
\begin{equation}
x'=x-vt,
\qquad
t'=t.
\end{equation}
That is Galilean invariance. The essential difference is that in the Galilean limit all observers share the same time. In special relativity they do not.

Once the notation is in place, the field Euler--Lagrange equation can also be written in a clean relativistic form:
\begin{equation}
\partial_\mu\left(\frac{\partial\mathcal L}{\partial(\partial_\mu\phi)}\right)
=
\frac{\partial\mathcal L}{\partial\phi}.
\end{equation}
This is exactly the same equation we were using before, but now the contraction of upper and lower indices displays its relativistic structure immediately.

\section{Massive Scalar Fields and the Klein--Gordon Equation}

The lecture now returns to dynamics with the sentence, in effect, \emph{let's talk about waves}. But before solving the wave equation, Susskind modifies the Lagrangian slightly. The extra term is as simple as possible: a quadratic term in the field itself. This is presented not as an abstract possibility but as the field-theory analog of the harmonic oscillator.

\begin{figure}
\centering
\includegraphics[width=0.82\textwidth]{lecture_05_figure_05.png}
\end{figure}

This board frame preserves exactly that pedagogical pairing. The upper line contains the field-theory Lagrangian; the lower line contains the ordinary oscillator comparison. The handwriting at the far right may look like \(M^2\phi^2\), but the lecture's spoken derivation clearly uses \(\mu^2\), and we standardize the notation that way:
\begin{equation}
\mathcal L
=
\frac12
\left[
\left(\frac{\partial\phi}{\partial t}\right)^2
-\left(\frac{\partial\phi}{\partial x}\right)^2
-\left(\frac{\partial\phi}{\partial y}\right)^2
-\left(\frac{\partial\phi}{\partial z}\right)^2
-\mu^2\phi^2
\right].
\end{equation}
The lower comparison line is just the oscillator Lagrangian
\begin{equation}
L_{\mathrm{osc}}=\frac{\dot\phi^2}{2}-\frac{\mu^2\phi^2}{2}.
\end{equation}
The field theory is the oscillator, now spread through space and supplemented by gradient terms.

Varying this Lagrangian is straightforward. The derivative with respect to \(\partial_t\phi\) gives \(\partial_t\phi\), so one more time derivative yields \(\partial_t^2\phi\). The spatial derivative terms similarly produce the negative Laplacian. Finally,
\begin{equation}
\frac{\partial\mathcal L}{\partial\phi}=-\mu^2\phi.
\end{equation}
Thus the equation of motion is
\begin{equation}
\partial_t^2\phi-\nabla^2\phi+\mu^2\phi=0.
\end{equation}
This is the Klein--Gordon equation.

A student then asks, naturally, where the equation came from. The answer is important for the tone of the chapter: it was not conjured out of folklore. It came from \emph{this} Lagrangian by the Euler--Lagrange rule. The lecture also remarks that this is essentially the unique simple invariant Lagrangian that leads to a linear equation of this type.

To solve the equation, we look for plane waves:
\begin{equation}
\phi(X)=e^{-i\omega t+i(k_xx+k_yy+k_zz)}.
\end{equation}
Once the notation is available, the lecture compresses this to
\begin{equation}
\phi(X)=e^{ik_\mu x^\mu},
\qquad
k_\mu=(-\omega,k_x,k_y,k_z).
\end{equation}
Differentiating twice gives
\begin{equation}
\partial_t^2\phi=-\omega^2\phi,
\qquad
\partial_x^2\phi=-k_x^2\phi,
\qquad
\partial_y^2\phi=-k_y^2\phi,
\qquad
\partial_z^2\phi=-k_z^2\phi.
\end{equation}
Substituting into the Klein--Gordon equation, we get
\begin{equation}
\left(-\omega^2+k_x^2+k_y^2+k_z^2+\mu^2\right)\phi=0.
\end{equation}
A nonzero solution therefore requires
\begin{equation}
\omega^2=k_x^2+k_y^2+k_z^2+\mu^2.
\end{equation}
Later in the lecture the same role is momentarily played by \(m^2\); we keep \(\mu^2\) here to match the Lagrangian from which the equation was derived.

The lecture immediately interprets this as the relativistic energy-momentum relation in field-theory language. There are two branches, \(\omega=\pm\sqrt{\mathbf k^2+\mu^2}\), and the general solution is a superposition of such plane waves. This is why the equation is easy to solve in principle: it is linear, and every solution can be decomposed into modes.

\subsection{Question \& Answer}

Why are complex exponentials allowed if the field is supposed to be real?

Because the complex exponential is only a computational convenience. By Euler's formula,
\begin{equation}
e^{i(\omega t-kx)}
=
\cos(\omega t-kx)+i\sin(\omega t-kx).
\end{equation}
The equation is linear, so if the complex exponential is a solution, then its real and imaginary parts are separately solutions. The cosine and sine are simply two real plane waves \(90^\circ\) out of phase. We use exponentials because differentiating them is easy, not because the field itself must be complex.

\section{Summary and Physical Interpretation}

At the end of the lecture Susskind returns one last time to the particle action. The coupling we wrote down earlier,
\begin{equation}
\mathcal L_{\mathrm{particle}}
=
-(m+g\phi)\sqrt{1-\dot x^{\,2}},
\end{equation}
has an immediate interpretation when the scalar field sits at a constant background value. In that case the particle behaves as though it has an effective mass
\begin{equation}
m_{\mathrm{eff}}=m+g\phi.
\end{equation}
A nonzero background field shifts the mass. If the bare mass were absent to begin with, a nonzero background could create a mass dynamically. This is the limited but important sense in which the lecture invokes the Higgs mechanism.

The point is then sharpened. Different particles may couple with different values of \(g\), and therefore the same background field can generate different masses for different particles. That is the physical meaning of the late Higgs remark. The lecture also briefly touches superconductivity and the possibility of a field giving an effective mass to the electromagnetic field, but it explicitly postpones that story until electrodynamics has been developed.

What remains from the lecture is a very clean chain of ideas. A spacetime picture yields the Doppler shift. A toy two-coordinate system explains action and reaction in the Lagrangian language. A scalar field enters the relativistic particle action, and in the nonrelativistic limit the coupling becomes ordinary potential energy. The same coupling then reappears as a source term in the field equation. Once the component formulas become unwieldy, relativistic notation is introduced to make invariance manifest. Finally a quadratic term in the field produces the Klein--Gordon equation and its plane-wave solutions. The mathematics never gets very ornate, but the method is already the modern one: write an invariant action, vary it, and let both particles and fields emerge from the same principle.